% from ega4-21b-fr.tex
\oldpage[IV-4]{293}
de $\mathscr{Z}^1_X\otimes_{\bb{Z}}\bb{Q}$ au-dessus de $X$ sont les \emph{cycles $1$-codimensionnels à coefficients rationnels}.
Comme les fibres de $\mathscr{Z}^1_X$ sont des $\bb{Z}$-modules sans torsion \sref{IV.21.6.3}, l'homomorphisme canonique $\mathscr{Z}^1_X\to\mathscr{Z}^1_X\otimes_{\bb{Z}}\bb{Q}$ est \emph{injectif}, de sorte que les cycles $1$-codimensionnels s'identifient à des cycles $1$-codimensionnels à coefficients rationnels.
\end{env}

\begin{env}[21.10.10]
\label{IV.21.10.10}
Nous allons voir que sous certaines conditions, on peut élargir la définition de $f^*(Z)$ donnée dans \sref{IV.21.10.3} pour un cycle $1$-codimensionnel $Z$ sur $X$, mais à condition de prendre pour $f^*(Z)$ un cycle $1$-codimensionnel à coefficients rationnels sur $X'$.
Le cas plus général où nous nous plaçons est celui où $f$ transforme tout point maximal de $X'$ en un point maximal de $X$, et où, en tout point $x'\in X'^{(1)}$, on a l'une des conditions (i), (ii), (iii) de \sref{IV.21.10.1} \emph{ou une quatrième condition} (en posant $x=f(x')$):
\begin{enumerate}
  \item[(iv)] $x\in X^{(1)}$, $\mathfrak{m}_x\notin\operatorname{Ass}(\sh{O}_{X,x})$, et \emph{en outre, si l'on pose} $A=\widehat{\sh{O}}_{X,x}$, $A'=\widehat{\sh{O}}_{X',x'}$, \emph{et si $K$ désigne l'anneau total des fractions de $A$, alors $A'$ est une $A$-algèbre finie et $K'=A'\otimes_AK$ est un $K$-module} libre.
\end{enumerate}

Notons alors $r_{x'}$ le rang du $K$-module libre $K'$, $q_{x'}$ le degré de $\boldsymbol{k}(x')=\boldsymbol{k}(A')$ sur $\boldsymbol{k}(x)=\boldsymbol{k}(A)$, et posons $\mu_{x'}=r_{x'}/q_{x'}$.
Pour un cycle $1$-codimensionnel de support contenu dans $T$, donné par \sref{IV.21.10.1.1} et tout $x'\in X'^{(1)}$, on définit $c_{x'}\in\bb{Q}$, comme égal au nombre $n_{x'}\in\bb{Z}$ lorsque l'on est dans l'un des cas $1\textsuperscript{o}$, $2\textsuperscript{o}$, $3\textsuperscript{o}$ de \sref{IV.21.10.1}; mais il reste ici une quatrième possibilité:
\begin{enumerate}
  \item[$4\textsuperscript{o}$] si $x'$ vérifie la condition (iv) ci-dessus, on prend $c_{x'}=\mu_{x'}n_x\in\bb{Q}$.
\end{enumerate}
\end{env}

\begin{env}[21.10.11]
\EGAReaderTarget{ega4.c04.statement.21-10-11}{IV.21.10.11}
Il faut encore prouver que lorsque la condition $4\textsuperscript{o}$ de \sref{IV.21.10.10} est vérifiée \emph{en même temps} qu'une des conditions $1\textsuperscript{o}$, $2\textsuperscript{o}$, $3\textsuperscript{o}$ de \sref{IV.21.10.1}, on a $n_{x'}=c_{x'}$.
C'est évident si $x\notin T$ puisque alors $n_x=0$.
Pour étudier les deux autres cas, notons que $\mathfrak{m}_x\sh{O}_{X',x'}$ est fermé pour la topologie $\mathfrak{m}_{x'}$-préadique, donc le complété de $\sh{O}_{X',x'}/\mathfrak{m}_x\sh{O}_{X',x'}$ pour cette dernière topologie est $A'/\mathfrak{m}_xA'$.
Si l'on est à la fois dans le cas $2\textsuperscript{o}$ et le cas $4\textsuperscript{o}$, $\sh{O}_{X',x'}/\mathfrak{m}_x\sh{O}_{X',x'}$ est discret pour la topologie $\mathfrak{m}_{x'}$-préadique, donc isomorphe à $A'/\mathfrak{m}_xA'$.
Comme $A'$ est une $A$-algèbre finie et \emph{plate} $(\mathbf{0}_{\mathrm{III}},10.2.3)$, c'est un $A$-module \emph{libre} $(\mathbf{0}_{\mathrm{III}},10.1.3)$, et le rang de $A'/\mathfrak{m}_xA'$ sur $A/\mathfrak{m}_xA=\boldsymbol{k}(x)$ est égal au rang de $A'$ sur $A$, donc aussi à celui de $K'$ sur $K$.
D'autre part ce rang est aussi égal au produit de la longueur $\lambda_{x'}$ de $\sh{O}_{X',x'}/\mathfrak{m}_x\sh{O}_{X',x'}$ (en tant que $\sh{O}_{X',x'}$-module, ou de $\boldsymbol{k}(x')$-module) par le rang $[\boldsymbol{k}(x'): \boldsymbol{k}(x)]=q_{x'}$, ce qui prouve la relation $\mu_{x'}=r_{x'}/q_{x'}=\lambda_{x'}$ dans ce cas.

Supposons enfin que l'on soit à la fois dans le cas $3\textsuperscript{o}$ et le cas $4\textsuperscript{o}$.
Alors, puisque $\dim(\sh{O}_{X,x})=1$, $\sh{O}_{X,x}$ est un anneau de valuation discrète, donc régulier.
D'autre part, $\sh{O}_{X',x'}$ est de dimension $1$, et puisque $\mathfrak{m}_{x'}\notin\operatorname{Ass}(\sh{O}_{X',x'})$, $\sh{O}_{X',x'}$ est un anneau de Cohen--Macaulay $(\mathbf{0},16.4.6)$;
enfin, puisque $A'$ est un $A$-module de type fini, $A'/\mathfrak{m}_xA'$ est un $\boldsymbol{k}(x)$-espace vectoriel de rang fini;
puisque $\sh{O}_{X',x'}/\mathfrak{m}_x\sh{O}_{X',x'}$ est contenu dans $A'/\mathfrak{m}_xA'$, c'est aussi un $\boldsymbol{k}(x)$-espace vectoriel de rang fini, donc un anneau artinien.
L'application de \sref{IV.6.1.5} montre alors que $\sh{O}_{X',x'}$ est un $\sh{O}_{X,x}$-module \emph{plat}, donc on est aussi dans le cas $2\textsuperscript{o}$, et on conclut par ce qui a été vu plus haut.

Cela étant, dans le cas que l'on a considéré, on posera
\[
  f^*(Z)=\sum_{x'\in X'^{(1)}}c_{x'}\cdot\overline{\{x'\}}.
\]
